\section{谷山--志村猜想（模性定理）的陈述}
\label{sec:ch10-ts}
% ============================================================================

\providecommand{\Zl}{\mathbb{Z}_\ell}
\providecommand{\Ql}{\mathbb{Q}_\ell}

终章的真正主角不是 Frey 曲线，
而是“每条有理椭圆曲线都来自模形式”这一句宏大断言。
从今天回头看，
我们会把它称为模性定理；
但在 20 世纪后半叶，
它曾经以“谷山--志村猜想”乃至“谷山--志村--韦伊猜想”的名字出现，
长期被视为解析数论、代数几何与 Galois 表示之间最神秘的桥。

\textbf{我们遇到了什么问题？}
到第\ref{ch:eichler-shimura}章时，
我们已经知道：
若 $f\in \Sk_2(\GammaO(N))$ 是归一化新形式，
则存在与之对应的模 Abel 簇 $A_f$，而当系数域 $K_f=\Q$ 时它就是一条椭圆曲线。
这说明“从模形式到椭圆曲线”是可能的。
但费马大定理真正需要的是反方向：
给定一条椭圆曲线 $E/\Q$，
它是否一定来自某个权 $2$ 新形式？
如果这个问题得不到肯定回答，
那么 Frey 曲线即使再古怪，
我们也无法把它送回模形式世界中制造矛盾。

\textbf{核心思想是什么？}
谷山与志村的洞见，是把“椭圆曲线”理解为模曲线的像。
级别 $N$ 的模曲线 $X_0(N)$ 参数化带有循环子群的椭圆曲线；
若某条椭圆曲线 $E$ 真与级别 $N$ 的模形式有关，
那么它应当不仅在 $L$-函数层面上与某个新形式匹配，
还应当接受一条来自 $X_0(N)$ 的满射有理映射。
这就是模参数化。
随着 Eichler--Shimura 与 Deligne 表示理论的发展，
这一猜想逐渐显现出四种彼此等价的外观：
它可以说成 $L$-函数一致、可以说成 Galois 表示一致、
可以说成存在模参数化，也可以说成 Hecke 特征值与 Frobenius 迹一致。

\textbf{引入之后能得到什么？}
一旦半稳定情形的猜想被证明，
Frey 曲线立刻落入其适用范围，
因为它正是一条半稳定椭圆曲线。
于是费马方程的假想解会自动给出一个权 $2$ 新形式；
而这正是 Ribet 水平下降可以继续处理的输入数据。
因此，Wiles 所证明的并不是一条与 FLT 外在相连的定理，
而是 FLT 证明链中不可缺少的中央齿轮。

\begin{figure}[ht]
\centering
\begin{tikzpicture}[>=Stealth, font=\small]
  \draw[very thick, blue!50!black] (0,0) -- (13.8,0);
  \foreach \x/\y in {0.6/1955,2.5/1960s,4.3/1967,6.2/1985,8.0/1986,9.5/1987,11.4/1995,13.0/2001}
    \draw[fill=white] (\x,0) circle (2.2pt) node[below=6pt] {\y};

  \node[above=8pt] at (0.6,0) {谷山提出模性问题};
  \node[above=8pt] at (2.5,0) {志村精确化猜想};
  \node[above=8pt] at (4.3,0) {Weil 提供几何证据};
  \node[above=8pt] at (6.2,0) {Frey 提出反常曲线};
  \node[above=8pt] at (8.0,0) {Ribet 证明水平下降};
  \node[above=8pt] at (9.5,0) {Serre 提出 $\varepsilon$ 猜想};
  \node[above=8pt] at (11.4,0) {Wiles 与 Taylor--Wiles 解决半稳定情形};
  \node[above=8pt] at (13.0,0) {BCDT 完成一般模性};
\end{tikzpicture}
\caption{谷山--志村猜想走向模性定理的时间线：从 1955 年的直觉性问题，到 1995 年的半稳定情形，再到 2001 年的一般情形。}
\label{fig:ch10-history-timeline}
\end{figure}


先给出终章真正要使用的版本。

\begin{theorem}[Wiles 的半稳定模性定理]
\label{thm:wiles}
设 $E/\Q$ 是半稳定椭圆曲线，导子为 $N$。
则存在权 $2$ 的归一化特征新形式
\[
  f\in \Sk_2(\GammaO(N))
\]
使得
\[
  L(s,E)=L(s,f).
\]
等价地，对每个素数 $\ell$ 都有
\[
  \rho_{E,\ell}\cong \rho_{f,\ell},
\]
并且存在模参数化
\[
  \phi\colon X_0(N)\twoheadrightarrow E.
\]
\end{theorem}

\begin{proof}[证明说明]
这一定理由 Wiles 在 1994 年首先证明，
并在与 Taylor 合作修补后于 1995 年完成发表。
真正的证明占据了后两节要介绍的全部技术：
Langlands--Tunnell 定理提供残余模性起点，
Mazur 形变理论与 Hecke 代数给出两个看似不同的局部对象 $R$ 与 $T$，
Taylor--Wiles 方法则证明它们同构。
本节不证明定理，
而是把它当作终章后续所有论证的主定理。
\end{proof}

\textbf{注。}
2001 年 Breuil--Conrad--Diamond--Taylor 把上面“半稳定”三字去掉，
从而证明了所有有理椭圆曲线都模。
不过对费马大定理来说，
半稳定情形已经完全足够，
因为 Frey 曲线恰好半稳定。

下面解释模性的四种等价表述。

\begin{proposition}[模性的四种面貌]
\label{prop:ch10-modularity-equivalent}
对一条有理椭圆曲线 $E/\Q$ 与一个权 $2$ 归一化新形式 $f\in \Sk_2(\GammaO(N))$，以下说法彼此等价，或至少在本章语境下可互相视为同一件事的不同表述：
\begin{itemize}
  \item[(i)] $L(s,E)=L(s,f)$；
  \item[(ii)] 对所有素数 $\ell$，有 $\rho_{E,\ell}\cong \rho_{f,\ell}$；
  \item[(iii)] 存在满射有理映射 $\phi\colon X_0(N)\twoheadrightarrow E$；
  \item[(iv)] 对几乎所有好素数 $p\nmid N$，有
  \[
    a_p(E)=a_p(f),
  \]
  其中 $a_p(E)=p+1-\#E(\F_p)$。
\end{itemize}
\end{proposition}

\begin{proof}[说明]
(i) 与 (iv) 只是 Euler 因子逐项一致的两种写法；
(ii) 与 (iv) 由第\ref{ch:eichler-shimura}章、第\ref{ch:galois}章的 Frobenius 迹公式联结起来；
(iii) 与 (i) 则由 Jacobian 上的 Hecke 特征部分与 Eichler--Shimura 关系联系。
在现代语言里，
这些表述由 Eichler--Shimura、Deligne、Faltings 同种定理与模参数化理论拼成一个闭环。
因此本章谈“$E$ 是模的”时，
读者可以在这四种面貌之间自由切换。
\end{proof}

历史上，这个猜想经历了一个漫长而戏剧化的成长过程。
1955 年谷山丰在日光会议上提出问题；
1960 年代志村五郎把它精确化，
Weil 则补充了理论证据并把它推向国际舞台；
1985 年 Frey 指出 FLT 可能是它的推论；
1986 年 Ribet 证明水平下降，
把“Frey 曲线若模，则 FLT 成立”变成严格定理；
1995 年 Wiles 与 Taylor--Wiles 证明半稳定情形；
2001 年 BCDT 完成一般情形。
这条时间线之所以动人，
是因为一个最初近乎“物理直觉”的猜想，
最后被锻造成了一个可完全落地的表示论定理。

\begin{example}[级别 $11$ 的模参数化]
\label{ex:ch10-x011-parametrization}
第4章与第6章已经说明
\[
  g(X_0(11))=\dim \Sk_2(\GammaO(11))=1,
\]
因此 $J_0(11)$ 是一维 Abel 簇，也就是一条椭圆曲线。
再由第7章的显式模型，
可取
\[
  E_{11}\colon y^2+y=x^3-x^2-10x-20.
\]
因为尖点 $[\infty]\in X_0(11)(\Q)$ 有理，
Abel--Jacobi 映射
\[
  \phi\colon X_0(11)\longrightarrow J_0(11),
  \qquad
  P\longmapsto [P]-[\infty]
\]
是一个定义在 $\Q$ 上的非平凡态射。
两边都是亏格 $1$ 的曲线，
且 $\phi([\infty])=0$，
故它实际上给出一个同构
\[
  X_0(11)\xrightarrow{\ \sim\ }J_0(11)\cong E_{11}.
\]
这就是级别 $11$ 的模参数化。
它已经把“模曲线映到椭圆曲线”这件事写成了一个完全显式的映射。
\end{example}

\textbf{注。}
有时文献把“模参数化”写成从解析商 $\Gamma_0(N)\backslash\HH^*$ 到 $E(\C)$ 的映射，
有时写成从代数曲线 $X_0(N)$ 到 $E$ 的有理态射。
这两种写法在本章中不必区分，
因为我们已经在第7章说明了模曲线的解析与代数模型是一回事。

这一节最值得牢牢记住的句子只有一句：
\emph{Frey 曲线若存在，它就必须模。}
这句简单的话，
实际上压缩了 Wiles 证明的全部重量。
下一节开始，我们只讨论这个重量是如何被组织起来的。
